\section{Теоретические сведения}

\subsection{Метод медленно меняющихся амплитуд}

В работе исследуется LC генератор с контуром в цепи сетки \cref{fig:th:1}, относящийся к разряду квазисинусоидальных автоколебательных систем.

\begin{figure}[H]
	\centering
	\includegraphics[width=0.4\linewidth]{th_1}
	\caption{Схема лабораторной установки}
	\label{fig:th:1}
\end{figure}

\begin{figure}[H]
	\centering
	\includegraphics[width=0.35\linewidth]{th_2}
	\caption{}
	\label{fig:th:2}
\end{figure}

Отправной точкой теоретического исследования этой системы являются уравнения Кирхгофа, которые, с учётом обозначений, принятых на \cref{fig:th:1}, запишутся в следующем виде:
\begin{align} \label{eq:th:1}
	L\dv{i}{t} + Ri + U = M\dv{i_a}{t} - E_0, \nonumber \\
	i = C\dv{U}{t} + i_c
\end{align}
где $i_a$ и $i_c$ анодные и сеточные токи лампы, зависящие в общем случае от анодного и сеточного напряжений. В дальнейшем будем учитывать их зависимость только от сеточных напряжений, т.е. пренебрежем реакцией анода. Качественный вид анодно-сеточной и сеточной характеристик приведен на \cref{fig:th:2}. В проводимом рассмотрении эти нелинейные функции будут аппроксимироваться полиномами.

Исключив из системы \cref{eq:th:1} ток и введя новые переменные
\begin{align} \label{eq:th:2}
	x = U + E_0, \qquad \tau = \frac{t}{\sqrt{LC}},
\end{align}
запишем исходную систему \cref{eq:th:1} в виде одного уравнения второго порядка
\begin{align} \label{eq:th:3}
	\dv[2]{x}{\tau} + x = \dv{}{\tau}\left[-\delta x + \sigma f_1(x)\right] - f_2(x).
\end{align}
Здесь $\sigma = \frac{M}{\sqrt{LC}}$, $\alpha = \frac{L}{M}$, $\delta = \frac{R\sqrt{C}}{\sqrt{L}}$, $f_1 = (i_a - a i_c)$, $f_2 = R i_c$. Заметим, что функция $f_1(x)$ включает в себя разность анодного и сеточного токов, но т.к. сеточный ток много меньше анодного, то не будет большой ошибкой считать $f_1 = i_a$.

Согласно \cref{eq:th:3}, в случае малых линейных потерь и нелинейностей рассматриваемая система будет близка к линейной и консервативной. К такой системе может быть применён метод медленно меняющихся амплитуд [2]. Следуя этому методу, решение \cref{eq:th:3} будем искать в виде
\begin{align} \label{eq:th:4}
	x = z e^{j\tau} + z^* e^{-j\tau}, \nonumber \\
	\dot{x} = j\left(z e^{j\tau} - z^* e^{-j\tau}\right).
\end{align}
Заметим, что второе равенство в \cref{eq:th:4} предполагает выполнение условия $\dot{z}e^{j\tau} + \dot{z}^*e^{-j\tau} = 0$. Подставляя \cref{eq:th:4} в уравнение \cref{eq:th:3} и усредняя результат подстановки по периоду, получим, с учётом медленности изменения комплексной амплитуды $Z$, укороченные уравнения
\begin{align} \label{eq:th:5}
	2j\dot{z} = -j\delta z + j\sigma\chi_1(z, z^*) - \chi_2(z, z^*),
\end{align}
где
\begin{align} \label{eq:th:6}
	\chi_1(z, z^*) = \frac{1}{2\pi}\int_{-\pi}^{\pi} f_1 e^{-j\tau}dt, \nonumber \\
	\chi_2(z, z^*) = \frac{1}{2\pi}\int_{-\pi}^{\pi} f_2 e^{-j\tau}dt
\end{align}
Аппроксимируя $f_1$ и $f_2$ степенными рядами
\begin{align} \label{eq:th:7}
	f_1(x) = \sum_{n=0}^{\infty}\frac{a_n x^n}{n!}, \qquad f_2(x) = \sum_{n=0}^{\infty}\frac{b_n x^n}{n!}
\end{align}
найдём для $\chi_1$ и $\chi_2$ выражения вида
\begin{align} \label{eq:th:8}
	\chi_1(z, z^*) = z\overline{f_1}(|z|^2), \qquad \chi_2(z, z^*) = z\overline{f_2}(|z|^2)
\end{align}
где
\begin{align} \label{eq:th:9}
	\overline{f_1}(|z|^2) = \sum_{m=0}^{\infty}\frac{a_{2m+1}(|z|^2)^m}{m!(m+1)!}, \qquad \overline{f_2}(|z|^2) = \sum_{m=0}^{\infty}\frac{b_{2m+1}(|z|^2)^m}{m!(m+1)!}
\end{align}
Полагая в уравнении \cref{eq:th:5} $z = \rho e^{j\phi}$, $z^* = \rho e^{-j\phi}$, перейдем от комплексных переменных $z$, $z^*$ к действительным модулю $\rho$ и фазе $\phi$:
\begin{align} \label{eq:th:10}
	\dot{\rho} = -\frac{1}{2}\left(\delta - \sigma\overline{f_1}(\rho^2)\right)\rho, \nonumber \\
	\dot{\phi} = \frac{1}{2}\overline{f_2}(\rho^2).
\end{align}
Из укороченных уравнений \cref{eq:th:10} видно, что условия возбуждения генератора и амплитуда установившихся автоколебаний определяются видом функции $\overline{f_1}$, т.е. в основном анодно-сеточной характеристикой лампы; поправка к частоте определяется функцией $\overline{f_2}$, т.е. сеточной характеристикой. Для отыскания установившихся (стационарных) значений амплитуды автоколебаний достаточно найти устойчивые состояния равновесия $\overline{\rho}$ первого из уравнений системы \cref{eq:th:10}. При этом из второго уравнения этой системы найдём поправку к частоте автоколебаний
\begin{align} \label{eq:th:11}
	\Delta\omega = \overline{f_2}(\overline{\rho}^2)/2,
\end{align}
задающую величину, на которую частота колебаний генератора будет отличаться от частоты колебательного контура. Тем самым, в режиме установившихся автоколебаний будем иметь \cref{eq:th:12,eq:th:13}
\begin{align} \label{eq:th:12}
	\phi = \frac{\overline{f_2}(\overline{\rho}^2)}{2}\tau + \phi_0
\end{align}
\begin{align} \label{eq:th:13}
	x = 2\overline{\rho}\cos(\tau + \phi) = 2\overline{\rho}\cos\left[\left(1 + \frac{\overline{f_2}(\overline{\rho}^2)}{2}\right)\tau + \phi_0\right].
\end{align}

\subsection{Стационарные режимы работы генератора [1]}

Прежде чем переходить к изучению различных режимов работы лампового генератора, выясним, при каких условиях справедлива та или иная аппроксимация нелинейной характеристики лампы. Как видно из \cref{eq:th:9}, на динамику генератора и стационарные режимы его работы влияют только нечетные члены степенного ряда нелинейной характеристики $f_1(x)$. Выделим из нелинейной функции нечетную ее часть с помощью соотношения \cref{eq:th:14}:
\begin{align} \label{eq:th:14}
	f_\text{н}(x) = \frac{f(x) - f(-x)}{2}.
\end{align}
Вид нечетной части характеристики, а, следовательно, и возможная её аппроксимация зависят от выбора рабочей точки, положение которой определяется постоянным смещением $E_0$ на сетке лампы. Рассмотрим возможные аппроксимации нелинейной функции $f_1$ и, как следствие, различные режимы работы генератора.

1. Напряжение смещения $E_0$ на управляющей сетке лампы (рабочая точка) выбрано так, что нечетная часть анодно-сеточной характеристики имеет вид, приведенный на \cref{fig:th:3}. Это возможно в том случае, когда напряжение смещения задано в точке максимальной крутизны характеристики. При этом функцию $f_1$ достаточно точно можно представить в виде полинома третьей степени, причем аппроксимация будет справедлива для части кривой, обозначенной сплошной линией на \cref{fig:th:3}:
\begin{align} \label{eq:th:15}
	f_1(x) = a_1 + \frac{a_2}{2}x^2 - \frac{a_3}{6}x^3.
\end{align}
Здесь и в дальнейшем все коэффициенты аппроксимирующего полинома будем считать положительными. Для аппроксимации \cref{eq:th:15} средняя крутизна анодно-сеточной характеристики, т.е. функция $\overline{f_1}$ примет вид:
\begin{align} \label{eq:th:16}
	\overline{f_1}(\rho^2) = a_1 - \frac{a_3}{2}\rho^2.
\end{align}
В этом случае закон изменения амплитуды колебаний, согласно системе \cref{eq:th:10}, будет определяться следующим уравнением:
\begin{align} \label{eq:th:17}
	\dot{\rho} = -\rho\overline{F}(\rho),
\end{align}
где $\overline{F}(\rho) = \left[\delta - \sigma\left(a_1 - \frac{a_3}{2}\rho^2\right)\right]/2$. Приравнивая правую часть этого уравнения нулю, получим стационарные значения амплитуд колебаний:
\begin{align} \label{eq:th:18}
	\overline{\rho_1} = 0, \qquad \overline{\rho_2} = \sqrt{\frac{2(-\delta + \sigma a_1)}{\sigma a_3}}.
\end{align}
Их устойчивость определяется корнями следующего характеристического уравнения:
\begin{align} \label{eq:th:19}
	p = \dv{}{\rho}\left[-\rho\overline{F}(\rho)\right]_{\rho=\overline{\rho}} = -\overline{F}(\overline{\rho}) - \overline{\rho}\overline{F}'(\overline{\rho}).
\end{align}

\begin{figure}[H]
	\centering
	\includegraphics[width=0.3\linewidth]{th_3}
	\caption{}
	\label{fig:th:3}
\end{figure}

Нетрудно видеть, что состояние системы с нулевым значением $\rho$ (невозбужденный генератор) будет устойчивым при $\overline{F}(0) > 0$, т.е. при $\sigma < \delta/a_1$ и неустойчивым в противном случае. Второе стационарное состояние устойчиво, если $\overline{F}'(\rho_2) = \sigma a_3\overline{\rho_2} > 0$, что выполняется при выбранных знаках коэффициентов. Выражение \cref{eq:th:19} определяет зависимость амплитуды колебаний от линейного декремента $\delta = R\sqrt{C/L}$, коэффициента связи $\sigma = M/\sqrt{LC}$, и коэффициентов аппроксимации анодно-сеточной характеристики $a_1$ и $a_3$. Можно построить бифуркационные диаграммы, выражающие зависимость $\overline{\rho}^2$ от любого из перечисленных параметров. Зависимость $\overline{\rho}^2$ от параметра $\sigma$ (величина обратной связи) приведена на \cref{fig:th:4a}, где точками обозначены устойчивые состояния, а крестиками неустойчивые. Из приведенной диаграммы видно, что схема возбуждается при условии $\sigma > \delta/a_1$ или, если перейти к параметрам схемы, при $M > RC/a_1$. В этом неравенстве величина $a_1$ крутизна анодно-сеточной характеристики в рабочей точке, иногда обозначаемая в литературе буквой $S$. Точка $\sigma = \delta/a_1$ на бифуркационной диаграмме называется точкой бифуркации здесь качественно меняется поведение системы. Фазовые портреты системы при $\sigma < \delta/a_1$ и $\sigma > \delta/a_1$ приведены на \cref{fig:th:4b}. Как видно из \cref{fig:th:4a,fig:th:4b}, предельный цикл устанавливается при сколь угодно малых начальных условиях, т.е. генератор обладает мягким режимом возбуждения по отношению к начальным условиям. Термин <<мягкий>> может употребляться и в другом смысле в качестве характеристики типа генератора, отличающегося плавным нарастанием установившегося значения амплитуды колебаний при медленном и непрерывном изменении параметра обратной связи. Приведенная на \cref{fig:th:4a} бифуркационная диаграмма соответствует именно такому типу генератора.

\begin{figure}[H]
	\centering
	\begin{minipage}{0.45\textwidth}
		\centering
		\includegraphics[width=0.9\linewidth]{th_4}
		\subcaption{Бифуркационная диаграмма}
		\label{fig:th:4a}
	\end{minipage}
	\begin{minipage}{0.45\textwidth}
		\centering
		\includegraphics[width=0.9\linewidth]{th_5}
		\subcaption{Фазовые портреты}
		\label{fig:th:4b}
	\end{minipage}
	\caption{}
	\label{fig:th:4}
\end{figure}

2. Рабочая точка выбрана так, что нечетная часть характеристики имеет вид, приведенный на \cref{fig:th:5}. Этому случаю чаще всего отвечают напряжения смещения, близкие к напряжению отсечки лампы. При этом характеристику лампы необходимо аппроксимировать полиномом не ниже пятой степени. Такой аппроксимации соответствует средняя крутизна
\begin{align} \label{eq:th:20}
	\overline{f_1}(\rho^2) = a_1 + \frac{a_3}{2}\rho^2 - \frac{a_5}{12}\rho^4
\end{align}
и, следовательно, (см. \cref{eq:th:10}) уравнение, определяющее стационарные значения амплитуды автоколебаний принимает вид
\begin{align} \label{eq:th:21}
	\rho\left[\delta - \sigma\left(a_1 + \frac{a_3}{2}\overline{\rho}^2 - \frac{a_5}{12}\overline{\rho}^4\right)\right] = 0.
\end{align}

\begin{figure}[H]
	\centering
	\includegraphics[width=0.3\linewidth]{th_6}
	\caption{}
	\label{fig:th:5}
\end{figure}

Бифуркационная диаграмма для этого случая приведена на \cref{fig:th:6}. Из него следует, что при значениях параметра $\sigma$, лежащих в интервале $\sigma' < \sigma < \sigma''$, возникновение автоколебаний возможно лишь, если начальное значение $\rho$ превышает соответствующее пороговое значение. В этом случае говорят о жёстком режиме возбуждения генератора. Бифуркационная диаграмма, изображённая на \cref{fig:th:6}, характеризуется гистерезисным эффектом установления и срыва автоколебаний при непрерывном изменении параметра $\sigma$. Системы с таким видом бифуркационной диаграммы относят к жёсткому типу генераторов. Для них, согласно \cref{fig:th:6}, мягкое возбуждение возможно лишь при $\sigma > \sigma''$, т.е. за пределами области гистерезиса. На \cref{fig:th:7} приведены фазовые портреты такого генератора для различных значений параметра $\sigma$.

\begin{figure}[H]
	\centering
	\includegraphics[width=0.4\linewidth]{th_7}
	\caption{}
	\label{fig:th:6}
\end{figure}

3. Если рабочая точка выбрана при положительных напряжениях на сетке лампы, то график нечетной части нелинейной характеристики имеет вид, приведенный на \cref{fig:th:8}. Здесь для аппроксимации анодно-сеточной характеристики лампы необходимо использовать полином седьмой степени. В этом случае функция $\overline{f_1}(\rho^2)$ и уравнение для определения ненулевых состояний равновесия запишутся в виде
\[
	\overline{f_1}(\rho^2) = a_1 + \frac{a_3}{2}\rho^2 - \frac{a_5}{12}\rho^4 - \frac{a_7}{144}\rho^6, \qquad \delta - \sigma\left(a_1 + \frac{a_3}{2}\overline{\rho}^2 - \frac{a_5}{12}\overline{\rho}^4 - \frac{a_7}{144}\overline{\rho}^6\right) = 0.
\]

\begin{figure}[H]
	\centering
	\includegraphics[width=0.85\linewidth]{th_8}
	\caption{}
	\label{fig:th:7}
\end{figure}
\begin{figure}[H]
	\centering
	\includegraphics[width=0.35\linewidth]{th_9}
	\caption{}
	\label{fig:th:8}
\end{figure}

При этом возможны два варианта бифуркационной диаграммы $\overline{\rho}^2(\sigma)$, изображённые на \cref{fig:th:9a,fig:th:9b}.

\begin{figure}[H]
	\centering
	\begin{minipage}{0.45\textwidth}
		\centering
		\includegraphics[width=0.9\linewidth]{th_10}
		\subcaption{Переход между предельными циклами}
		\label{fig:th:9a}
	\end{minipage}
	\begin{minipage}{0.45\textwidth}
		\centering
		\includegraphics[width=0.9\linewidth]{th_11}
		\subcaption{Срыв колебаний при $\sigma = \sigma_0$}
		\label{fig:th:9b}
	\end{minipage}
	\caption{}
	\label{fig:th:9}
\end{figure}

В варианте \cref{fig:th:9a} при уменьшении параметра $\sigma$ происходит скачкообразный переход с большего предельного цикла на меньший и при дальнейшем уменьшении $\sigma$ предельный цикл плавно исчезает. В варианте \cref{fig:th:9b} при уменьшении параметра $\sigma$ до $\sigma_0$ колебания срываются до нуля. Описанная разновидность автоколебательных систем относится к генераторам сложно-жесткого типа.
